\documentclass[10pt]{article}
\usepackage[margin=0.62in]{geometry}
\usepackage{booktabs}
\usepackage[hidelinks]{hyperref}
\usepackage{natbib}
\usepackage{amsmath}
\usepackage{amssymb}
\usepackage{array}
\usepackage{microtype}
\title{Multilinear Discriminant Analysis for Face Recognition:\\A Paper-Proximate ORL Study}
\author{COMP3314 Group 27}
\date{Fall 2026}
\begin{document}
\maketitle

\begin{abstract}
Yan et al. proposed Multilinear Discriminant Analysis (MDA), which learns interrelated discriminative projections along the modes of an image tensor. We studied the method on the original 112$\times$92 ORL faces, comparing Eigenfaces PCA, Fisherfaces LDA, MDA/2-1, and MDA/2-2 with Euclidean 1-nearest-neighbour. On 30 matched splits at fixed dimensions, MDA/2-1 performs best in our implementation; MDA/2-2 is highly sensitive to projection-axis scaling. Five-repeat searches over all 10,130 feasible MDA/2-2 dimension pairs show a large gap from the paper's selected scores. Our work follows the stated equations and update order, but exact reproduction is not possible from the published details: preprocessing defaults, split seeds, eigenvector scaling, singular-scatter handling, and the dimension-selection protocol are not fully specified.
\end{abstract}

\section{Problem and method}
Vector-based PCA and LDA flatten an image into one long feature vector. MDA instead keeps an image as a matrix and learns projections along its row and column modes. For image $X_i\in\mathbb{R}^{h\times w}$, the two-mode feature is $U^\top X_i V$. MDA's discriminant tensor criterion extends Fisher's between- versus within-class idea:
\begin{equation*}
J(U,V)=
\frac{\sum_c n_c\left\|U^\top(\bar X_c-\bar X)V\right\|_F^2}
{\sum_c\sum_{i:y_i=c}\left\|U^\top(X_i-\bar X_c)V\right\|_F^2}.
\end{equation*}
Here $\bar X_c$ and $\bar X$ are the gallery class and overall means, and $n_c$ is the gallery count for identity $c$. With one mode fixed, the other is updated from class-weighted between-class scatter and pooled within-class scatter, then a generalized eigenproblem is solved. MDA/2-1 updates one mode; MDA/2-2 alternates row then column updates from identity initialization.

\begin{center}
\fbox{\parbox{0.94\linewidth}{\centering
\textbf{ORL face matrix} $X$
\quad$\longrightarrow$\quad
\textbf{Learn row and column axes} $U^\top X V$
\quad$\longrightarrow$\quad
\textbf{Euclidean 1-NN identity}\\[3pt]
\small MDA/2-1 learns one axis; MDA/2-2 alternates both axes.}}
\end{center}

This connects to course topics: PCA is unsupervised variance reduction; LDA uses labels to separate classes; each MDA mode update also solves a generalized eigenproblem. The classifier uses Euclidean distances in the resulting feature space, so changing axis scales can change predictions.

\section{Data and protocol}
We use the Cambridge ORL archive: 400 grayscale images, 40 identities, 10 images per identity, at 112$\times$92 pixels. Each image is equalized independently. The paper specifies histogram equalization but not its software or output levels; we use 64 output levels as a MATLAB-default-aligned approximation. Per identity, 2, 3, 4, or 5 images form the gallery; the rest are probes. All projections are fit on gallery images only, and every method shares each deterministic split. The seed rule is $3314+1000g+r$, where $g$ is gallery size and $r$ is zero-based repeat index.

The fixed-setting study uses 30 splits per gallery size: PCA with 36 components; Fisherfaces with full $N-C$ PCA pre-reduction and up to 36 LDA components; MDA/2-1 with six row components; and MDA/2-2 with $6\times6$ components. Euclidean 1-nearest-neighbour is used throughout. The primary MDA results use unit-L2 eigenvector axes; within-scatter normalization is retained as a sensitivity setting. For the paper-style MDA stopping rule, after iteration 2 each projection change is compared with $m'_k m_k\epsilon$. The article does not report $\epsilon$; we use $10^{-4}$ and a 10-iteration cap. A relative $10^{-6}$ ridge regularizes singular within-class scatter.

\section{Results}
\begin{table}[h]
\centering
\small
\caption{Fixed-setting probe accuracy, mean $\pm$ sample SD over 30 matched splits (percent). Paper values are the reported best MDA/2-2 scores, selected over feature dimensions; they are not directly comparable to our fixed $6\times6$ results.}
\begin{tabular}{crrrrrr}
\toprule
Gallery & PCA & Fisherfaces & MDA/2-1 & MDA/2-2 & MDA/2-2 & Paper MDA/2-2\\
images/person & & & & unit-L2 & within-scatter & best\\
\midrule
2 & 73.9 $\pm$ 2.5 & 73.1 $\pm$ 2.8 & 81.4 $\pm$ 3.0 & 34.8 $\pm$ 5.3 & 5.4 $\pm$ 1.9 & 90.1\\
3 & 84.1 $\pm$ 1.8 & 84.9 $\pm$ 2.9 & 89.9 $\pm$ 2.2 & 52.6 $\pm$ 5.6 & 23.9 $\pm$ 10.3 & 95.0\\
4 & 88.2 $\pm$ 3.0 & 88.0 $\pm$ 2.4 & 92.0 $\pm$ 2.2 & 75.8 $\pm$ 3.6 & 64.0 $\pm$ 6.0 & 98.3\\
5 & 91.6 $\pm$ 2.0 & 90.4 $\pm$ 2.2 & 94.7 $\pm$ 1.7 & 87.5 $\pm$ 2.8 & 81.4 $\pm$ 3.8 & 99.0\\
\bottomrule
\end{tabular}
\end{table}

MDA/2-1 leads our fixed-setting comparison at all four gallery sizes. MDA/2-2 improves as the gallery grows, but its accuracy changes sharply with axis normalization. This is not a controlled reconstruction of the paper's chosen dimensions, and the paper does not identify its exact eigenvector scaling convention.

\begin{table}[h]
\centering
\small
\caption{Highest mean accuracy for a single MDA/2-2 dimension pair across five splits, percent. The selected pair is chosen using those same probe scores, so these maxima are optimistic diagnostics, not unbiased test estimates.}
\begin{tabular}{crrr}
\toprule
Gallery images/person & Within-scatter best & Unit-L2 best & Unit-L2 pair\\
\midrule
2 & 62.44 & 61.81 & $112\times86$\\
3 & 57.29 & 82.29 & $64\times46$\\
4 & 77.92 & 90.75 & $40\times26$\\
5 & 88.60 & 94.90 & $30\times22$\\
\bottomrule
\end{tabular}
\end{table}

The gallery-size-3 dip occurs only in the probe-selected within-scatter grid. In a separate 30-split nested-gallery diagnostic with dimension fixed at $6\times6$, adding a third image improves mean accuracy from 34.57\% to 53.05\% with unit-L2 axes, and from 5.47\% to 26.70\% with within-scatter axes. Thus the selected-grid dip does not show that adding a gallery image generally hurts recognition. The full unit-L2 and within-scatter heatmaps and split-level CSVs are linked from the repository overview.

\section{Critical analysis}
The paper's MDA/2-2 Table I values are 90.1\%, 95.0\%, 98.3\%, and 99.0\% for gallery sizes 2--5 \citep{Yan2007MDA}. Our selected grid values remain lower, particularly at gallery sizes 2--4. The gap cannot be attributed to one factor from the available evidence. The paper does not give random split seeds, full dimension ranges, exact preprocessing software/defaults, the eigenvector scaling convention, the treatment of singular scatter, or its dimension-selection protocol. Our gallery-only fitting avoids probe leakage in the learned projections, but selecting the best dimension by probe accuracy reuses the reported labels and inflates the grid maxima. A less biased estimate needs training-only dimension selection followed by an untouched outer test set; that would be a distinct evaluation protocol.

\section{Implementation, efficiency, and reproducibility}
The original full grids remain the canonical experimental tables: each normalization has five repeats over 10,130 feasible pairs, or 202,600 MDA/2-2 fits. They were generated with the direct implementation. The optimized sweep is now the default for future dimension searches. It reuses class grouping and generalized eigensystem factorizations, then forms each requested-size projection with the same shape-specific matrix multiply as the direct solver. A cache that reused truncated full bases caused rare iteration and accuracy differences on nearly singular gallery-size-2 fits; that path was corrected. The 17 configurations flagged by the interrupted run were rechecked under both normalizations: corrected direct and cached fits had identical predictions, accuracies, iteration counts, projections, distances, and objectives. Recomputed direct fits also matched the canonical accuracy and iteration counts. No exhaustive rerun was scientifically necessary.

On the M3 Max, the corrected cached path measured 47.48 fits/s at 14 workers and one Accelerate thread (0.91 GiB peak RSS) over 60 representative fits; the direct path measured 40.50 fits/s (1.19 GiB). The projected time for 202,600 fits is about 71 minutes from this representative batch, not a completed full-run timing. Classification was below 1\% of fit time; scatter construction and eigendecomposition were the largest profiled stages. The interrupted partial checkpoint is retained for audit and is not used in the canonical tables. No remote HPC access was assumed or used.

The paper provides a foundation for MDA's tensor formulation and alternating update method, but this study does not numerically reproduce its best-tuned table. The strongest result is a transparent ORL implementation with preserved full-grid evidence, explicit assumptions, and a tested efficient path for future sweeps.

\section*{Reproducibility and attribution}
Source, environment, raw split-level results, dimension summaries, visualizations, validation records, and run commands are in the private Group 27 repository. ORL data are loaded from the Cambridge archive and are not redistributed. Cite the original paper \citep{Yan2007MDA} and the ORL database. The local overview links directly to both full grids, the corrected fixed-setting results, and the optimization validation.

\bibliographystyle{plainnat}
\bibliography{references}
\end{document}
